\section{Monoidal reachability for acyclic cluster algebras}
\label{Sec : monoidal reachability for acyclic cluster algebras}

\subsection{Cluster categories of acyclic quivers}
\label{ssec : cluster categories of acyclic quivers}

Let $Q$ be a finite acyclic quiver and $\C Q$ be the path algebra.
Let $\operatorname{mod}\C Q$ be the category of finitely generated
modules over $\C Q$ and $D^b(\operatorname{mod}\C Q)$ be the bounded
derived category of $\operatorname{mod}\C Q$.
Let $\tau$ be the Auslander--Reiten translation in
$D^b(\operatorname{mod}\C Q)$ and $\Sigma$ be the shift functor.

\begin{Def}[{\cite[\S1]{BMRRT}}]
\label{Def : cluster category}
We define the \emph{cluster category}
$\mathcal{C}_Q \coloneq
D^b(\operatorname{mod}\C Q)/(\tau\Sigma^{-1})$ as follows:
\[
\operatorname{obj}(\mathcal{C}_Q)
= \operatorname{obj}(D^b(\operatorname{mod}\C Q)),
\]
\[
\operatorname{Hom}_{\mathcal{C}_Q}(X, Y)
= \bigoplus_{i \in \Z}
\operatorname{Hom}_{D^b(\operatorname{mod}\C Q)}
\bigl((\tau\Sigma^{-1})^iX, Y\bigr).
\]
\end{Def}

\begin{Prop}[{\cite{Kel05, BMRRT}}]
\label{Prop : properties of cluster categories}
Let $Q$ be an acyclic quiver.
\begin{enumerate}
    \item[(i)] $\mathcal{C}_Q$ is a triangulated category.
    \item[(ii)] The canonical functor
    $\pi \colon D^b(\operatorname{mod}\C Q) \longrightarrow
    \mathcal{C}_Q$ is a triangle functor.
    \item[(iii)] $\mathcal{C}_Q$ is Krull--Schmidt.
    \item[(iv)]
    \[
    \operatorname{ind}\mathcal{C}_Q
    = \operatorname{ind}(\operatorname{mod}\C Q)
    \sqcup \{\Sigma P_i \mid i \in Q_0\},
    \]
    where $\operatorname{ind}(-)$ is the set of isomorphism classes
    of indecomposable objects and $P_i$ is the indecomposable
    projective module associated to $i \in Q_0$.
\end{enumerate}
\end{Prop}

\begin{Def}[{\cite{BMRRT}}]
\label{Def : rigid and cluster tilting objects}
We say that $X \in \mathcal{C}_Q$ is \emph{rigid} if
$\operatorname{Ext}^1_{\mathcal{C}_Q}(X, X) = 0$.
We say that $T \in \mathcal{C}_Q$ is a \emph{cluster-tilting object}
if $T$ is rigid and $T$ has a maximal number of non-isomorphic
indecomposable direct summands.
An object $X \in \mathcal{C}_Q$ is called \emph{basic}
if all of its indecomposable direct summands are non-isomorphic.
\end{Def}

\begin{Thm}[{\cite[Theorem 3.3, Theorem 5.1]{BMRRT}}]
\label{Thm : cluster tilting summands and mutation}
\begin{enumerate}
    \item[(i)] Any cluster-tilting object has $|Q_0|$
    non-isomorphic indecomposable direct summands.
    \item[(ii)] For a basic cluster-tilting object
    $T = \bigoplus_{i \in Q_0}T_i$ and $k \in Q_0$,
    there exists a unique rigid indecomposable object
    $T_k^* \not\simeq T_k$ such that
    \[
    T_k^* \oplus \bigoplus_{i \in Q_0 \setminus \{k\}}T_i
    \]
    is a basic cluster-tilting object.
\end{enumerate}
\end{Thm}

For a basic cluster-tilting object
$T = \bigoplus_{i \in Q_0}T_i$ and $k \in Q_0$, we call
$T_k^* \oplus \bigoplus_{i \in Q_0 \setminus \{k\}}T_i$
the \emph{mutation} of $T$ at $k$.

\begin{Thm}[{\cite[Proposition 3.5]{BMRRT}}]
\label{Thm : additive reachability}
Let $Q$ be an acyclic quiver.
Then any basic cluster-tilting object $T$ is reachable,
i.e., $T$ is obtained from the basic cluster-tilting object
$\bigoplus_{i \in Q_0}P_i$ by finitely many mutations.
\end{Thm}

For $\mathbf{a} \in \Z_{\geq 0}^{Q_0}$, we set
\[
P(\mathbf{a}) \coloneq
\bigoplus_{i \in Q_0}P_i^{\oplus a_i}.
\]

\begin{Prop}[{\cite{Palu08}}]
\label{Prop : presentations and index}
For $X \in \mathcal{C}_Q$, there exists a triangle
\[
P(\mathbf{g}_1) \longrightarrow P(\mathbf{g}_0)
\longrightarrow X \xrightarrow{+1}
\]
for $\mathbf{g}_0, \mathbf{g}_1 \in \Z_{\geq 0}^{Q_0}$,
and $\mathbf{g}_0 - \mathbf{g}_1 \in \Z^{Q_0}$ is independent
of the choice of the presentation.
\end{Prop}

For $X \in \mathcal{C}_Q$, we define
$\operatorname{ind}X \coloneq
\mathbf{g}_0 - \mathbf{g}_1 \in \Z^{Q_0}$.

\begin{Def}[{\cite[\S2.2]{CK06}}]
\label{Def : quiver Grassmannians}
Let $M \in \operatorname{mod}\C Q$
and $\mathbf{d} \in \Z_{\geq 0}^{Q_0}$.
We define the \emph{quiver Grassmannian}
$\operatorname{Gr}_{\mathbf{d}}(M)$ by
\[
\{N \subset M \mid N \text{ is a submodule of } M
\text{ such that } \underline{\dim}N = \mathbf{d}\}.
\]
\end{Def}

\begin{Def}[{\cite{CK06}}]
\label{Def : cluster character}
Let $X \in \mathcal{C}_Q$.
We write $\Sigma X = M \oplus \Sigma P$
for $M \in \operatorname{mod}\C Q$ and a projective module $P$.
We define the \emph{cluster character}
$\mathrm{CC}(X) \in \Z[x_i^{\pm 1} \mid i \in Q_0]$ by
\[
\mathrm{CC}(X) \coloneq
\mathbf{x}^{\operatorname{ind}X}
\sum_{\mathbf{d} \in \Z_{\geq 0}^{Q_0}}
\chi\bigl(\operatorname{Gr}_{\mathbf{d}}(M)\bigr)
\mathbf{x}^{B_{Q^{\mathrm{op}}}\mathbf{d}},
\]
where $\chi$ is the Euler characteristic and $B_{Q^{\mathrm{op}}}$ is
the skew-symmetric matrix corresponding to $Q^{\mathrm{op}}$.
\end{Def}

\begin{Thm}[{\cite[Theorem 4]{CK06}}]
\label{Thm : cluster tilting objects and clusters}
Let $\overline{\mathcal{A}}_{Q^{\mathrm{op}}}$ be the cluster algebra with no frozen variables 
with the initial seed
\[
(\mathbf{x}_{t_0}, B_{t_0})
= \bigl((x_i)_{i \in Q_0}, B_{Q^{\mathrm{op}}}\bigr).
\]
Then
\[
\mathrm{CC} \colon \mathcal{C}_Q \longrightarrow
\Z[x_i^{\pm 1} \mid i \in Q_0]
\]
induces the bijection
\[
\{\text{indecomposable rigid objects in } \mathcal{C}_Q\}
\longrightarrow
\{\text{cluster variables of }
\overline{\mathcal{A}}_{Q^{\mathrm{op}}}\}.
\]
Moreover, it induces the bijection
\[
\{\text{basic cluster-tilting objects in } \mathcal{C}_Q\}
\longrightarrow
\{\text{clusters of }\overline{\mathcal{A}}_{Q^{\mathrm{op}}}\},
\]
and it is compatible with the mutations.
\end{Thm}

Note that for $X \in \mathcal{C}_Q$ such that
$\operatorname{ind}X = \mathbf{g}$, we have
\[
\deg^{t_0}\mathrm{CC}(X) = \mathbf{g}
\quad\text{in }\overline{\mathcal{A}}_{Q^{\mathrm{op}}}.
\]


\subsection{$F$-polynomials and canonical decompositions}
\label{ssec : F polynomials and canonical decompositions}

Let $Q$ be an acyclic quiver.
For $\mathbf{g} \in \Z^{Q_0}$, we set
\[
\operatorname{PHom}(\mathbf{g})
\coloneq \operatorname{Hom}_{\C Q}
\bigl(P([-\mathbf{g}]_+), P([\mathbf{g}]_+)\bigr).
\]
For $d \in \operatorname{PHom}(\mathbf{g})$, set
$X_d \coloneq \operatorname{Cone}(d)$.
Note that $\operatorname{ind}X_d = \mathbf{g}$
for $d \in \operatorname{PHom}(\mathbf{g})$.

\begin{Def}[{\cite[Definition 3.3]{DF15}}]
\label{Def : E invariant of presentations}
Let $\mathbf{g}_i \in \Z^{Q_0}$ and
$d_i \in \operatorname{PHom}(\mathbf{g}_i)$ for $i = 1, 2$.
We define
\[
E(d_1, d_2)
\coloneq
\operatorname{Hom}_{K^b(\operatorname{proj}\text{-}\C Q)}
(d_1, \Sigma d_2),
\]
where $K^b(\operatorname{proj}\text{-}\C Q)$ is the bounded homotopy
category of $\operatorname{proj}\text{-}\C Q$.
We set
\[
e(\mathbf{g}_1, \mathbf{g}_2)
\coloneq
\min\left\{
\dim E(d_1, d_2)
\;\middle|\;
 d_i \in \operatorname{PHom}(\mathbf{g}_i) \text{ for } i = 1, 2
\right\}.
\]
We define
\[
\operatorname{ext}(\mathbf{g}_1, \mathbf{g}_2)
\coloneq
\min\left\{
\dim\operatorname{Ext}^1_{\mathcal{C}_Q}(X_{d_1}, X_{d_2})
\;\middle|\;
 d_i \in \operatorname{PHom}(\mathbf{g}_i) \text{ for } i = 1, 2
\right\}.
\]
Since
\[
\dim E,\ \dim\operatorname{Ext}^1_{\mathcal{C}_Q}
\colon
\operatorname{PHom}(\mathbf{g}_1) \times \operatorname{PHom}(\mathbf{g}_2)
\longrightarrow \Z
\]
are upper semicontinuous \cite[Corollary 3.7]{DF15},
$e$ and $\operatorname{ext}$ are the generic values of
$\dim E$ and $\dim\operatorname{Ext}^1_{\mathcal{C}_Q}$.
\end{Def}

% P_1 \longrightarrow P_2^{\oplus(n-1)}
% P_1 \longrightarrow P_2^{\oplus(n-1)}
\begin{Prop}[{\cite[Lemma 3.3]{Palu08}}]
\label{Prop : generic cluster extensions and E invariants}
Let $\mathbf{g}_1, \mathbf{g}_2 \in \Z^{Q_0}$.
Then
\[
\operatorname{ext}(\mathbf{g}_1, \mathbf{g}_2)
= e(\mathbf{g}_1, \mathbf{g}_2) + e(\mathbf{g}_2, \mathbf{g}_1).
\]
\end{Prop}

\begin{Def}[{\cite[Definition 4.3]{DF15}}]
\label{Def : canonical decomposition of indices}
Let $\mathbf{g} \in \Z^{Q_0}$ and $d \in \operatorname{PHom}(\mathbf{g})$.
We say that $d$ is \emph{indecomposable} if $d$ is nonzero and
is not decomposable into a direct sum of nontrivial presentations.
We say that $\mathbf{g}$ is \emph{indecomposable} if a generic
$d \in \operatorname{PHom}(\mathbf{g})$ is indecomposable.
For $\mathbf{g} \in \Z^{Q_0}$,
$\mathbf{g} = \mathbf{g}_1 \oplus \cdots \oplus \mathbf{g}_s$ is a
\emph{canonical decomposition} if each $\mathbf{g}_i$ is indecomposable and
a generic element of $\operatorname{PHom}(\mathbf{g})$ has the form
$f_1 \oplus \cdots \oplus f_s$ with
$f_i \in \operatorname{PHom}(\mathbf{g}_i)$.
\end{Def}

\begin{Thm}[{\cite[Theorem 4.4]{DF15}}]
\label{Thm : characterization of canonical decompositions}
Let $\mathbf{g} \in \Z^{Q_0}$.
Then $\mathbf{g} = \mathbf{g}_1 \oplus \cdots \oplus \mathbf{g}_s$
is the canonical decomposition
if and only if each $\mathbf{g}_i$ is indecomposable, and
$e(\mathbf{g}_i, \mathbf{g}_j) = 0$ whenever $i \neq j$.
\end{Thm}

\begin{Prop}[{\cite[Corollary 21]{DW02}}]
\label{Prop : number of real and imaginary canonical summands}
Let $\{\mathbf{g}_1, \ldots, \mathbf{g}_r,
\mathbf{g}_{r+1}, \ldots, \mathbf{g}_{r+s}\}$ be a subset
of $\Z^{Q_0}$ satisfying the following conditions:
\begin{itemize}
    \item
    $\mathbf{g}_1 \oplus \cdots \oplus \mathbf{g}_{r+s}$
    is the canonical decomposition
    of $\sum_{i=1}^{r+s}\mathbf{g}_i$,
    \item
    each $\mathbf{g}_i$ $(1 \leq i \leq r)$ is an index of some
    indecomposable rigid object,
    \item
    each $\mathbf{g}_i$ $(r+1 \leq i \leq r+s)$ is not an index of any
    indecomposable rigid object.
\end{itemize}
Then we have $r + 2s \leq |Q_0|$.
\end{Prop}

\begin{Cor}
\label{Cor : full canonical decomposition and cluster tilting}
Let $\{\mathbf{g}_i \mid 1 \leq i \leq |Q_0|\}$ be a subset of $\Z^{Q_0}$
such that $\mathbf{g}_1 \oplus \cdots \oplus \mathbf{g}_{|Q_0|}$
is the canonical
decomposition of $\sum_{i=1}^{|Q_0|}\mathbf{g}_i$.
Then there exists a cluster-tilting object
$T = \bigoplus_{i=1}^{|Q_0|}T_i$ such that
$\operatorname{ind}T_i = \mathbf{g}_i$ for any $1 \leq i \leq |Q_0|$.
\end{Cor}

\begin{proof}
By Proposition~\ref{Prop : number of real and imaginary canonical summands},
each $\mathbf{g}_i$ is the index of some indecomposable rigid object $T_i$.
Since $T_i$ is rigid, $T_i$ is a generic object in
$\{X_{d_i} \mid d_i \in \operatorname{PHom}(\mathbf{g}_i)\}$.
Therefore,
\[
0 = \operatorname{ext}(\mathbf{g}_i, \mathbf{g}_j)
= \dim\operatorname{Ext}^1_{\mathcal{C}_Q}(T_i, T_j).
\]
Thus $T = \bigoplus_{i=1}^{|Q_0|}T_i$ is rigid and has $|Q_0|$
nonisomorphic direct summands, so $T$ is a basic cluster-tilting object.
\end{proof}

For $\Sigma X = M \oplus \Sigma P \in \mathcal{C}_Q$, we put
\[
F_X(z_i \mid i \in Q_0)
\coloneq
\sum_{d \in \Z_{\geq 0}^{Q_0}}
\chi\bigl(\operatorname{Gr}_d(M)\bigr)\mathbf{z}^d
\in \Z[z_i \mid i \in Q_0]
\]
and call it the \emph{$F$-polynomial} of $X$.

\begin{Prop}[{\cite[\S 3.4]{Plam13}}]
\label{Prop : generic cluster characters and F polynomials}
Let $\mathbf{g} \in \Z^{Q_0}$.
Then the map
\[
\operatorname{PHom}(\mathbf{g})
\longrightarrow \Z[x_i^{\pm 1} \mid i \in Q_0],
\qquad d \longmapsto \mathrm{CC}(X_d)
\]
has a generic value.
Moreover, the map
\[
\operatorname{PHom}(\mathbf{g})
\longrightarrow \Z[z_i \mid i \in Q_0],
\qquad d \longmapsto F_{X_d}(z_i \mid i \in Q_0)
\]
has a generic value.
We write it
$F_{\mathbf{g}}^{Q,\mathrm{gen}}(z_i \mid i \in Q_0)
\coloneq F_{\mathbf{g}}^{\mathrm{gen}}(z_i \mid i \in Q_0)$.
\end{Prop}

\begin{Prop}[{\cite{Schofield92}, \cite[Corollary 3.26]{Fei23}}]
\label{Prop : generic tropical F polynomials and E invariants}
For $\mathbf{g}_1, \mathbf{g}_2 \in \Z^{Q_0}$, we have
\[
F_{\mathbf{g}_1}^{\mathrm{gen}}[\mathbf{g}_2] = e(\mathbf{g}_1, \mathbf{g}_2).
\]
\end{Prop}

\begin{Rem}
\label{Rem : comparison with Fei notation}
Note that, in the notation of \cite{Fei23},
if $M = \operatorname{Ker}(\check{\delta})$, then
$f_M = F_{-\check{\delta}}^{\mathrm{gen}}[-]$.
\end{Rem}

\begin{Def}
\label{Def : Newton polytope}
For
\[
f(z_i \mid 1 \leq i \leq m)
= \sum_{\mathbf{v} \in \Z^m}c_{\mathbf{v}}\mathbf{z}^{\mathbf{v}}
\in \Z[z_i^{\pm 1} \mid 1 \leq i \leq m],
\]
the \emph{Newton polytope} of $f$ is the convex hull of
$\{\mathbf{v} \in \Z^m \mid c_{\mathbf{v}} \neq 0\}$ in $\mathbb{R}^m$.
We write the Newton polytope of $f$ as $\operatorname{Newt}(f)$.
\end{Def}

\begin{Lem}
\label{Lem : vertices of Newton polytopes}
Let $0 \neq F(z_i \mid i \in Q_0)
\in \Z[z_i^{\pm 1} \mid i \in Q_0]$
and $p \in \mathbb{R}^{Q_0}$.
The following are equivalent:
\begin{enumerate}
    \item[(i)]
    $p$ is a vertex of $\operatorname{Newt}(F)$.
    \item[(ii)]
    There exist $\{\mathbf{g}_i \in \Z^{Q_0} \mid i \in Q_0\}$ such that
    they are linearly independent, and
    \[
    F\left[\sum_{i \in Q_0}\mathbf{g}_i\right]
    = \sum_{i \in Q_0}F[\mathbf{g}_i],
    \qquad
    \{p\}
    = \bigcap_{i \in Q_0}
    \left\{\gamma \in \mathbb{R}^{Q_0}
    \;\middle|\; \mathbf{g}_i^{\mathsf{T}}\gamma = F[\mathbf{g}_i]\right\}.
    \]
\end{enumerate}
\end{Lem}

\begin{proof}
Assume (i).
Then the normal cone
\[
N_F(p)
= \left\{\gamma \in \mathbb{R}^{Q_0}
\;\middle|\;
\gamma\cdot(p-\mathbf{v}) \geq 0
\text{ for all } \mathbf{v} \in \operatorname{Newt}(F)\right\}
\]
is a $|Q_0|$-dimensional cone.
Therefore, there exist
$\{\mathbf{g}_i \in \Z^{Q_0} \mid i \in Q_0\} \subset N_F(p)$
such that they are linearly independent.
Then we have
\[
\begin{aligned}
F\left[\sum_{i \in Q_0}\mathbf{g}_i\right]
&= \max\left\{
\left(\sum_i \mathbf{g}_i\right)^{\mathsf{T}}\mathbf{v}
\;\middle|\; \mathbf{v} \in \operatorname{Newt}(F)\right\} \\
&= \left(\sum_i \mathbf{g}_i\right)^{\mathsf{T}}p.
\end{aligned}
\]
Therefore, we have
\[
F\left[\sum_i \mathbf{g}_i\right] = \sum_i F[\mathbf{g}_i],
\qquad
\{p\} = \bigcap_{i \in Q_0}
\left\{\gamma \in \mathbb{R}^{Q_0}
\;\middle|\; \mathbf{g}_i^{\mathsf{T}}\gamma = F[\mathbf{g}_i]\right\}.
\]
Assume (ii).
Let $p_0 \in \operatorname{Newt}(F)$ be a point which satisfies
\[
F\left[\sum_i \mathbf{g}_i\right]
= \left(\sum_i \mathbf{g}_i\right)^{\mathsf{T}}p_0.
\]
Then we have
\[
\begin{aligned}
0
&= \sum_i F[\mathbf{g}_i] - F\left[\sum_i \mathbf{g}_i\right] \\
&= \sum_i F[\mathbf{g}_i]
- \left(\sum_i \mathbf{g}_i\right)^{\mathsf{T}}p_0.
\end{aligned}
\]
Since $F[\mathbf{g}_i] \geq \mathbf{g}_i^{\mathsf{T}}p_0$ for all $i \in Q_0$,
we have $F[\mathbf{g}_i] = \mathbf{g}_i^{\mathsf{T}}p_0$ for all $i$.
Therefore,
\[
p_0 \in \bigcap_{i \in Q_0}
\left\{\gamma \in \mathbb{R}^{Q_0}
\;\middle|\; \mathbf{g}_i^{\mathsf{T}}\gamma = F[\mathbf{g}_i]\right\}.
\]
Because
\[
\left|
\bigcap_{i \in Q_0}
\left\{\gamma \in \mathbb{R}^{Q_0}
\;\middle|\; \mathbf{g}_i^{\mathsf{T}}\gamma = F[\mathbf{g}_i]\right\}
\right| = 1
\]
by the assumption, we have $p = p_0$.
Since
\[
\operatorname{Newt}(F)
\subset
\bigcap_{i \in Q_0}
\left\{\gamma \in \mathbb{R}^{Q_0}
\;\middle|\; \mathbf{g}_i^{\mathsf{T}}\gamma \leq F[\mathbf{g}_i]\right\}
\]
and $p \in \operatorname{Newt}(F)$ is a vertex of
\[
\bigcap_{i \in Q_0}
\left\{\gamma \in \mathbb{R}^{Q_0}
\;\middle|\; \mathbf{g}_i^{\mathsf{T}}\gamma \leq F[\mathbf{g}_i]\right\},
\]
$p$ is a vertex of $\operatorname{Newt}(F)$.
\end{proof}

\begin{Prop}[{\cite[Lemma 5.2, Theorem 6.4]{Fei23}}]
\label{Prop : cluster cones and generic Newton polytopes}
Let $\mathbf{g} \in \Z^{Q_0}$.
For a cluster-tilting object $T = \bigoplus_{i \in Q_0}T_i$,
\[
F_{\mathbf{g}}^{\mathrm{gen}}\left[\sum_i c_i\operatorname{ind}T_i\right]
= \sum_i c_i F_{\mathbf{g}}^{\mathrm{gen}}[\operatorname{ind}T_i]
\]
holds for all $c_i \in \Z_{\geq 0}$.
Moreover, for any vertex $p$ of
$\operatorname{Newt}(F_{\mathbf{g}}^{Q,\mathrm{gen}})$,
there exists some cluster-tilting object
$T = \bigoplus_i T_i$ such that
\[
\{p\} = \bigcap_{i \in Q_0}
\left\{\gamma \in \mathbb{R}^{Q_0}
\;\middle|\;
\gamma^{\mathsf{T}}(\operatorname{ind}T_i)
= F_{\mathbf{g}}^{\mathrm{gen}}[\operatorname{ind}T_i]\right\}.
\]
\end{Prop}


\subsection{Monoidal reachability for acyclic cluster algebras}
\label{ssec : monoidal reachability for acyclic cluster algebras}

Let
$\iota \colon K(\mathscr{C})_{\Z}
\xrightarrow{\sim}
\overline{\mathcal{A}}(\mathcal{S}_{\mathbf{x}}^{\Lambda}),J = 
\{1, \ldots, N\},J_{\mathrm{uf}},J_{\mathrm{fr}}$
be as in Subsection~\ref{ssec : Lambda cluster algebra structure}.
In this subsection, assume that
$\overline{\mathcal{A}}(\mathcal{S}_{\mathbf{x}}^{\Lambda})$
is of acyclic type.
We fix $t_1 \in \mathbb{T}$ such that the principal part $B_{t_1}$
of the exchange matrix $\widetilde{B}_{t_1}$ corresponds to
an acyclic quiver $Q^{\mathrm{op}}$.
Note that $J_{\mathrm{uf}} = Q_0$.
Since $t_1$ is injective-reachable, the following facts hold by
\cite[Theorem 1.3]{Plam13} and
\cite[Theorem 5.2.2]{Qin24}.

\begin{Lem}
\label{Lem : generic basis and compatible pointedness}
The set
\[
\left\{
L_{\widehat{\mathbf{g}}}
\coloneq \mathbf{x}_{t_1}^{\widehat{\mathbf{g}}}
F^{\mathrm{gen}}_{\mathbf{g}}
(\widehat{y}_{i;t_1} \mid i \in J_{\mathrm{uf}})
\;\middle|\; \widehat{\mathbf{g}} \in \Z^J
\right\}
\]
forms a basis of $\mathcal{U}(\mathcal{S}_{\mathbf{x}}^{\Lambda})$,
where $\mathbf{g} = \operatorname{pr}_{\mathrm{uf}}
\widehat{\mathbf{g}} \in \Z^{J_{\mathrm{uf}}} = \Z^{Q_0}$.
Moreover, each $L_{\widehat{\mathbf{g}}}$ is compatibly pointed and
$\deg^{t_1}L_{\widehat{\mathbf{g}}} = \widehat{\mathbf{g}}$.
\end{Lem}

This basis is called the \emph{generic basis} of
$\mathcal{U}(\mathcal{S}_{\mathbf{x}}^{\Lambda})$.
Note that the generic basis is not necessarily positive.

\begin{Prop}
\label{Prop : tropical F dominates generic tropical F}
Let $L_1, L_2 \in \Irr\mathscr{C}$ and
\[
\widehat{\mathbf{g}}_i
\coloneq \deg^{t_1}\iota([L_i]) \in \Z^J
\qquad\text{for }i = 1, 2.
\]
Let $\mathbf{g}_i
\coloneq \operatorname{pr}_{\mathrm{uf}}\widehat{\mathbf{g}}_i$
for $i = 1, 2$.
Then we have
\[
F^{t_1}_{\iota([L_1])}[\mathbf{g}_2]
\geq F^{Q,\mathrm{gen}}_{\mathbf{g}_1}[\mathbf{g}_2].
\]
\end{Prop}

\begin{proof}
It is enough to prove that
\[
\operatorname{Newt}(F^{t_1}_{\iota([L_1])})
\supset \operatorname{Newt}(F^{Q,\mathrm{gen}}_{\mathbf{g}_1}).
\]
Let us prove that each vertex $p$ of
$\operatorname{Newt}(F^{Q,\mathrm{gen}}_{\mathbf{g}_1})$
is a vertex of $\operatorname{Newt}(F^{t_1}_{\iota([L_1])})$.
By Proposition~\ref{Prop : cluster cones and generic Newton polytopes},
there exists a basic cluster-tilting object
$T = \bigoplus_{i \in Q_0}T_i$ such that
\[
F^{\mathrm{gen}}_{\mathbf{g}_1}
\left[\sum_i c_i\operatorname{ind}T_i\right]
= \sum_i c_i F^{\mathrm{gen}}_{\mathbf{g}_1}
[\operatorname{ind}T_i]
\qquad\text{for all }c_i \in \Z_{\geq 0},
\]
and
\[
\{p\} = \bigcap_{i \in Q_0}
\left\{\gamma \in \mathbb{R}^{Q_0}
\;\middle|\;
\gamma^{\mathsf{T}}\operatorname{ind}T_i
= F^{\mathrm{gen}}_{\mathbf{g}_1}[\operatorname{ind}T_i]
\right\}.
\]
By Lemma~\ref{Lem : vertices of Newton polytopes},
it is enough to prove that
\[
F^{t_1}_{\iota([L_1])}
\left[\sum_i c_i\operatorname{ind}T_i\right]
= F^{\mathrm{gen}}_{\mathbf{g}_1}
\left[\sum_i c_i\operatorname{ind}T_i\right]
\qquad\text{for all }c_i \in \Z_{\geq 0}.
\]
By Theorems~\ref{Thm : additive reachability} and
\ref{Thm : cluster tilting objects and clusters},
there exists $w \in \mathbb{T}$ such that
\[
\operatorname{ind}T_i
= \operatorname{pr}_{\mathrm{uf}}\deg^{t_1}x_{i;w}
\qquad\text{for }i \in J_{\mathrm{uf}}.
\]
Let
$\{\widehat{\mathbf{g}}_i^t \mid t \in \mathbb{T}\}
\in S_Y(\Z^{\max})$
be such that
$\widehat{\mathbf{g}}_i^{t_1} = \widehat{\mathbf{g}}_i$
for $i = 1, 2$.
By Proposition~\ref{Prop : tropical invariants of compatibly pointed elements},
Proposition~\ref{Prop : compatible pointedness of simple modules}, and
Lemma~\ref{Lem : generic basis and compatible pointedness}, we have
\[
\begin{aligned}
&\widehat{\mathbf{g}}_1^{\mathsf{T}}\Lambda_{t_1}
\bigl(\deg^{t_1}\mathbf{x}_{w}^{\mathbf{c}}\bigr)
+ 2F^{t_1}_{\iota([L_1])}
\left[\sum_i c_i\operatorname{ind}T_i\right] \\
&\qquad = (\widehat{\mathbf{g}}_1^{w})^{\mathsf{T}}
\Lambda_{w}
\bigl(\deg^{w}\mathbf{x}_{w}^{\mathbf{c}}\bigr)
+ 2F^{w}_{\iota([L_1])}
\left[\sum_i c_i\operatorname{pr}_{\mathrm{uf}}
\deg^{w}x_{i;w}\right]
\end{aligned}
\]
and
\[
\begin{aligned}
&\widehat{\mathbf{g}}_1^{\mathsf{T}}\Lambda_{t_1}
\bigl(\deg^{t_1}\mathbf{x}_{w}^{\mathbf{c}}\bigr)
+ 2F^{\mathrm{gen}}_{\mathbf{g}_1}
\left[\sum_i c_i\operatorname{ind}T_i\right] \\
&\qquad = (\widehat{\mathbf{g}}_1^{w})^{\mathsf{T}}
\Lambda_{w}
\bigl(\deg^{w}\mathbf{x}_{w}^{\mathbf{c}}\bigr)
+ 2F^{w}_{L_{\widehat{\mathbf{g}}_1}}
\left[\sum_i c_i\operatorname{pr}_{\mathrm{uf}}
\deg^{w}x_{i;w}\right].
\end{aligned}
\]
Since
\[
\sum_{i \in J_{\mathrm{uf}}}c_i\operatorname{pr}_{\mathrm{uf}}
\deg^{w}x_{i;w}
= \sum_{i \in J_{\mathrm{uf}}}c_i e_i
\in \Z_{\geq 0}^{J_{\mathrm{uf}}},
\]
we have
\[
F^{w}_{\iota([L_1])}
\left[\sum_i c_i\operatorname{pr}_{\mathrm{uf}}
\deg^{w}x_{i;w}\right]
= F^{w}_{L_{\widehat{\mathbf{g}}_1}}
\left[\sum_i c_i\operatorname{pr}_{\mathrm{uf}}
\deg^{w}x_{i;w}\right]
\]
by Proposition~\ref{Prop : compatible pointedness and bipointedness}.
Then we have
\[
F^{t_1}_{\iota([L_1])}
\left[\sum_i c_i\operatorname{ind}T_i\right]
= F^{\mathrm{gen}}_{\mathbf{g}_1}
\left[\sum_i c_i\operatorname{ind}T_i\right]
\qquad\text{for all }c_i \in \Z_{\geq 0}.
\]
\end{proof}

\begin{Cor}
\label{Cor : d dominates generic extensions}
Let $L_1, L_2 \in \Irr\mathscr{C}$.
Let $\widehat{\mathbf{g}}_i
\coloneq \deg^{t_1}\iota([L_i])$ and
$\mathbf{g}_i
\coloneq \operatorname{pr}_{\mathrm{uf}}\widehat{\mathbf{g}}_i$.
Then we have
\[
\mathfrak{d}(L_1, L_2)
\geq e(\mathbf{g}_1, \mathbf{g}_2)
+ e(\mathbf{g}_2, \mathbf{g}_1).
\]
\end{Cor}

\begin{proof}
By Theorem~\ref{Thm : tropical F invariants and Lambda inequalities},
Proposition~\ref{Prop : tropical F dominates generic tropical F}, and
Proposition~\ref{Prop : generic tropical F polynomials and E invariants}, we
have
\[
\begin{aligned}
\mathfrak{d}(L_1, L_2)
&\geq F^{t_1}_{\iota([L_1])}[\mathbf{g}_2]
+ F^{t_1}_{\iota([L_2])}[\mathbf{g}_1] \\
&\geq F^{\mathrm{gen}}_{\mathbf{g}_1}[\mathbf{g}_2]
+ F^{\mathrm{gen}}_{\mathbf{g}_2}[\mathbf{g}_1] \\
&= e(\mathbf{g}_1, \mathbf{g}_2)
+ e(\mathbf{g}_2, \mathbf{g}_1).
\end{aligned}
\]
\end{proof}

\begin{Ex}
\label{Ex : positive self d invariant}
Let
\[
A = \begin{pmatrix}2 & -n \\ -n & 2\end{pmatrix}
\qquad\text{for }n \in \Z_{\geq 3}.
\]
Consider the corresponding quiver Hecke algebra $R$ and let
$\widetilde{w} \coloneq s_1s_2s_1s_2$.
Let $\mathscr{C} = \mathscr{C}_w$ and
\[
\iota \colon K(\mathscr{C})_{\Z}
\xrightarrow{\sim}
\overline{\mathcal{A}}(\mathcal{S}_{\mathbf{x}}^{\Lambda})
\]
be the monoidal categorification as in
Subsection~\ref{ssec : Lambda cluster algebra structure}.
Here, $J = \{1, 2, 3, 4\}$, $J_{\mathrm{uf}} = \{1, 2\}$,
and $J_{\mathrm{fr}} = \{3, 4\}$.
Let $t_1 = t_0$.
Then $Q^{\mathrm{op}} = (1 \xrightarrow{n} 2)$
and $Q = (1 \xleftarrow{n} 2)$.
Note that $\dim P_1 = (1, 0)$ and $\dim P_2 = (n, 1)$
for $P_1, P_2 \in \operatorname{mod}\C Q$.
Let $\mathbf{g} = (-1, n-1)^{\mathsf{T}}$.
We can identify
\[
\operatorname{PHom}(\mathbf{g})
= \operatorname{Hom}_{\C Q}(P_1, P_2^{\oplus(n-1)})
\simeq M_{(n-1)\times n}(\C).
\]
For $d_1, d_2 \in \operatorname{PHom}(\mathbf{g})$
corresponding to $A, B \in M_{(n-1)\times n}(\C)$,
\[
\dim E(d_1, d_2)
= \dim\operatorname{Hom}_{\C Q}(P_1, P_2^{\oplus(n-1)})
- \dim\operatorname{Im}\delta_{A,B},
\]
where
\[
\delta_{A,B} \colon
\C \times M_{(n-1)\times(n-1)}(\C)
\longrightarrow M_{(n-1)\times n}(\C);
\qquad (a, C) \longmapsto Ba - CA.
\]
Therefore
\[
\dim E(d_1, d_2)
\geq n(n-1)-\bigl(1+(n-1)^2\bigr)
= n-2.
\]
Since $\delta_{A,B}$ is injective for generic $A, B$, we have
$e(\mathbf{g}, \mathbf{g}) = n-2$.
Let $L \in \Irr\mathscr{C}$ be such that
$\operatorname{pr}_{\mathrm{uf}}\deg^{t_1}\iota([L]) = \mathbf{g}$.
For example, such an $L$ appears as a composition factor of
$M_{1;\mu_1t_1} \otimes M_{2;\mu_2t_1}$.
Then $\mathfrak{d}(L, L) \geq 2(n-2) > 0$
by Corollary~\ref{Cor : d dominates generic extensions}.
In particular, there is a simple module whose self $\mathfrak{d}$-invariant
is non-zero.
\end{Ex}

\begin{Thm}
\label{Thm : monoidal reachability for acyclic cluster algebras}
For $L \in \Irr\mathscr{C}$,
$L$ is real if and only if $L$ is reachable.
Moreover, $L$ is real prime if and only if
$L$ corresponds to a cluster variable.
\end{Thm}

\begin{proof}
Let $\mathcal{M} = \{L_i \mid i \in J\}$ be a maximal real set
such that $L$ is a monomial in $\mathcal{M}$.
Let $\widehat{\mathbf{g}}_i \coloneq \deg^{t_1}\iota([L_i])$
and $\mathbf{g}_i
\coloneq \operatorname{pr}_{\mathrm{uf}}\widehat{\mathbf{g}}_i$
for $i \in J$.
Note that $\widehat{\mathbf{g}}_i = e_i$ for $i \in J_{\mathrm{fr}}$.
Let $\{\delta_1, \ldots, \delta_s\} \subset \Z^{Q_0}$
be the set of indecomposable indices which appear in
the canonical decomposition of some $\mathbf{g}_i$
for $i \in J_{\mathrm{uf}}$.
By Corollary~\ref{Cor : d dominates generic extensions}, we have
$e(\delta_i, \delta_j) = 0$ for $1 \leq i, j \leq s$.
Since
$\{\mathbf{g}_i \mid i \in J_{\mathrm{uf}}\}
\subset \Z^{J_{\mathrm{uf}}}$
is linearly independent by
Theorem~\ref{Thm : Laurent families and maximal real sets} and
Proposition~\ref{Prop : degrees of Laurent families form an integral basis},
we obtain that $s \geq |Q_0|$.
By Corollary~\ref{Cor : full canonical decomposition and cluster tilting},
$s = |Q_0|$ and there exists a basic cluster-tilting object
$T = \bigoplus_{i \in Q_0}T_i$ such that
$\delta_i = \operatorname{ind}T_i$ for $i \in J_{\mathrm{uf}}$.
By Theorem~\ref{Thm : additive reachability},
there exist $w \in \mathbb{T}$ and a permutation $\sigma$
of $J_{\mathrm{uf}}$ such that
\[
\operatorname{pr}_{\mathrm{uf}}\deg^{t_1}x_{\sigma(i);w}
= \delta_i.
\]
Let $\mathbf{g}_i = \sum_{j \in J_{\mathrm{uf}}}a_{ij}\delta_j$
for $i \in J_{\mathrm{uf}}$ and $a_{ij} \in \Z_{\geq 0}$.
Since
\[
\operatorname{pr}_{\mathrm{uf}}\deg^{t_1}\iota([L_i])
= \operatorname{pr}_{\mathrm{uf}}\deg^{t_1}
\iota\left(\left[
\bigotimes_{k \in J_{\mathrm{uf}}}
M_{\sigma(k);w}^{\otimes a_{ik}}
\right]\right),
\]
there exist $\mathbf{a}, \mathbf{b}
\in \Z_{\geq 0}^{J_{\mathrm{fr}}}$ such that
\[
\begin{aligned}
&\deg^{t_1}\iota\left(\left[
L_i \otimes \bigotimes_{j \in J_{\mathrm{fr}}}
M_{j;t_1}^{\otimes a_j}\right]\right) \\
&\qquad = \deg^{t_1}\iota\left(\left[
\bigotimes_{k \in J_{\mathrm{uf}}}
M_{\sigma(k);w}^{\otimes a_{ik}}
\otimes \bigotimes_{j \in J_{\mathrm{fr}}}
M_{j;t_1}^{\otimes b_j}\right]\right).
\end{aligned}
\]
By Lemma~\ref{Lem : injectivity of degrees for Laurent families},
we have
\[
L_i \otimes \bigotimes_{j \in J_{\mathrm{fr}}}
M_{j;t_1}^{\otimes a_j}
\simeq \bigotimes_{k \in J_{\mathrm{uf}}}
M_{\sigma(k);w}^{\otimes a_{ik}}
\otimes \bigotimes_{j \in J_{\mathrm{fr}}}
M_{j;t_1}^{\otimes b_j}.
\]
By Corollary~\ref{Cor : unique prime tensor factorization in C},
we have
\[
L_i \simeq \bigotimes_{k \in J_{\mathrm{uf}}}
M_{\sigma(k);w}^{\otimes a_{ik}}.
\]
Therefore, $L$ is a monomial in a monoidal seed
$\{M_{i;w} \mid i \in J\}$.
In particular, if $L$ is real prime,
$L$ is isomorphic to some $M_{i;w}$ by primeness.
\end{proof}
